\chapter{Quality Control and Validation}
\label{ch:quality}

\section{Pipeline checks}
The saved final quality-control report marks the applicable pipeline checks as \pass. It records no failed, warning, or not-run checks. The main validation results are summarized in \cref{tab:qc-summary}.

\small
\begin{longtable}{@{}p{0.30\textwidth}p{0.61\textwidth}@{}}
\caption{Selected results from the saved final pipeline quality-control report.}
\label{tab:qc-summary}\\
\toprule
\textbf{Check} & \textbf{Recorded result}\\
\midrule
\endfirsthead
\caption[]{Selected quality-control results (continued).}\\
\toprule
\textbf{Check} & \textbf{Recorded result}\\
\midrule
\endhead
\rowcolors{2}{ProjectGrey}{white}
Split sizes & Train 10,000; validation 100; test 1,000.\\
Reference captions & 5 per selected image; 55,500 records total.\\
Processed image integrity & 11,100 decoded; no missing processed images.\\
Split separation & No image-ID overlap between any pair of splits.\\
Exact duplicates & 0 duplicate groups by SHA-256.\\
Perceptual-hash screen & 9 within-split candidates; 0 cross-split candidates at Hamming distance 6.\\
Image embeddings & Expected split counts; 512 dimensions; finite and L2-normalized; manifest order checked.\\
Caption embeddings & 50,000 train vectors; 512 dimensions; unique IDs and source order checked.\\
Caption-index leakage & 0 validation/test caption IDs in the training index.\\
FAISS indexes & Image index: 10,000 vectors; caption index: 50,000; index readback passed.\\
Image retrieval & 20 of 20 queries returned results; top-5; 0 self-matches.\\
\bottomrule
\end{longtable}

\smallskip
\textit{Scope note:} The final quality-control report records zero failed, warning, and not-run checks among the checks it implements. The separate requirement-level compliance audit remains PARTIAL.


The reported checks cover expected split sizes, five captions for each selected image, processed-image presence and decoding, manifest separation, exact and perceptual duplicate screening, embedding dimensions and values, vector-to-ID mappings, FAISS index readback, caption-index exclusion of held-out IDs, and retrieval results. Perceptual-hash screening found nine within-split candidates and no cross-split candidates at the inclusive Hamming-distance threshold of 6. A perceptual hash is a screening signal; it is not proof of duplicate identity.

\section{Pipeline QC is not assignment compliance}
The quality-control report's \code{phase1\_complete: true} refers to the checks implemented by that pipeline. A separate requirement-level audit, \code{reports/PHASE1\_FINAL\_COMPLIANCE\_REPORT.md}, rates overall compliance as \partialstatus. It identifies missing or incomplete evidence that is outside the QC gates, including the unresolved test-size requirement, unavailable assignment PDF, unrecorded historical EXIF/conversion counts, partial caption-mapping verification, and an incomplete unified environment snapshot. Therefore, a QC PASS must not be read as proof that every project or assignment requirement has been met.
